% Time-stamp: <2026-09-12 20:44:08 komori>
% Copyright (c) 2021 Yasushi Komori
% All rights reserved.


%%
%% use .bux instead of .aux, which might be created in nonbraille mode.
%%
% in \document
%
% \UseOneTimeHook{begindocument/before}% <- AtEndPreamble
% ...
% \begingroup
% \@input{\jobname.aux}%
% \endgroup
% \if@filesw
%   \immediate\openout\@mainaux\jobname.aux
%   \immediate\write\@mainaux{\relax}%
% \fi
% ...
% \UseOneTimeHook{begindocument}% <- AtBeginDocument



%
\AddToHook{begindocument/before}{%
  %% 
  %% mecab
  %%
  % \ifx\brl@mecab@prog\empty%
  %   \brl@atbegdoc_setmecabprog%
  % \fi%
  % %
  % \ifx\brl@mecab@sysdic\empty%
  %   \brl@atbegdoc_setmecabsysdir%
  % \fi%
  % % 
  % \ifbrl@option@file%
  %   \edef\brl@mecab@translator{%
  %     \brl@mecab@prog -p -Obraille-file \brl@mecab@sysdic \brl@mecab@userdic \jobname.mcb -o \jobname.mrt%
  %   }%
  % \else%
  %   \edef\brl@mecab@translator{%
  %     \brl@mecab@prog \brl@mecab@sysdic \brl@mecab@userdic -p -Obraille%
  %   }%
  % \fi%
  %
  \let\brl@orig@bibcite\bibcite%
  \let\brl@orig@newlabel\newlabel%
  \let\bibcite\@gobbletwo%
  \def\newlabel{\@newl@bel{brl@@r}}% if aux exists, include it with \brl@XXX instead of \r@XXX
  \let\ifbrl@@filesw\if@filesw% save the original switch
   %
  \@fileswfalse% invalidate openout \@mainaux with suffix .aux in \document
}

\AtBeginDocument{%
  \let\bibcite\brl@orig@bibcite%
  \let\newlabel\brl@orig@newlabel%
  \begingroup%\@floatplacement\@dblfloatplacement
    \makeatletter\let\@writefile\@gobbletwo%
    \global\let\@multiplelabels\relax%
    \@input{\jobname.bux}% use bux as aux for braille.
  \endgroup%
  \let\if@filesw\ifbrl@@filesw% restore the original switch
%  \@gobble\fi%
  \if@filesw%
    \immediate\openout\@mainaux\jobname.bux%
    \immediate\write\@mainaux{\relax}%
  \fi%
  %
  %
  % set script for print
  \ifbrl@option@hires%
    \ifdefined\brl@bse@outputdevice%
      \brl@bse@outputlanguage%
    \else%
      \PackageError{jmbraille}{The serial port device for your printer is not specified. Give the device by \string\BrailleSetOutputDevice}{}%
    \fi%
    \brl@dbc_init% dbc
  \fi%
  %
  %
  \brl@plot_resource% register the resource for pdf
  %
  %
  %
  \ifbrl@lualatex%
    \pdfextension mapline {+DixBrailleA <[Braille.enc <<DixBraille-\brl@braille@fontname.otf}%
    \pdfextension mapline {+DixBrailleP <[Braille.enc <<DixBraille-Plain.otf}%
  \else%
    \AtBeginDvi{%
      \special{pdf:mapline DixBrailleA Braille.enc DixBraille-\brl@braille@fontname.otf}%
      \special{pdf:mapline DixBrailleP Braille.enc DixBraille-Plain.otf}%
    }%    
  \fi%
  % 
  % 
  \edef\brl@document@inputlinenob{\the\inputlineno}%
  \edef\brl@document@inputlinenoe{\the\inputlineno}%
  \edef\brl@document@basegrouplevel{\the\currentgrouplevel}%
  %
  \ifdefined\brl@document@nonbraillelinepage% if there is no nonbraille version.
    \let\brl@document_nonbraille@linepage\brl@document@nonbraillelinepage%
  \else%
    \let\brl@document_nonbraille_parsepage\empty%
  \fi%
  %
  \brl@orig@everypar{\brl@document}%
  %
  \ifbrl@math@altcolon%
    \brl@document_math_cmddef{rel}{\colon}{altcolon}%
    \brl@document_math_cmddef{colon}{:}{colon}%
  \fi%
  %
  \brl@ueb_pqm%
}%

\endinput


%%% Local Variables:
%%% mode: japanese-latex
%%% backup-each-save-backup : t
%%% TeX-master: t
%%% End:
